\section{Open questions}

\begin{enumerate}
  \item \textbf{Finite-temperature extension.} Does Trotterised PITE extend
        cleanly to preparing thermal Gibbs states $\rho = e^{-\beta H}/Z$
        (not just ground-state projection), and how does the per-Trotter-step
        ancilla success probability scale with $\beta$ and system size for a
        mixed target?
        \emph{Basis:} the paper's per-gadget $\alpha^2 = e^{2|c_k|\Delta\tau}$
        is derived for pure-state input; purification-based Gibbs preparation
        is not treated.
        \emph{Next steps:} implement PITE on a doubled (system + purifier)
        register for a 2--3-qubit toy Hamiltonian, initialise the purifier in
        the maximally entangled state, arrive at $\rho(\beta)$ by tracing out
        the purifier, and compare against exact $e^{-\beta H}$; measure
        cumulative success probability vs.\ $\beta$ at fixed spectral gap.

  \item \textbf{Tensor-network warm start.} Can PITE be warm-started from an
        MPS/PEPS approximation of the ground state so that the number of
        Trotter steps $r$ --- and hence $p \sim e^{-cr}$ --- shrinks by orders
        of magnitude for gapped 1D and quasi-1D Hamiltonians?
        \emph{Basis:} Fig.~7 starts from $|+\rangle^{\otimes n}$ with modest
        ground-state overlap; if $\epsilon = 1 - |\langle\psi_0|\psi_{\rm init}\rangle|^2$
        is small the required $\beta$ shrinks logarithmically in $\epsilon$
        and shot cost drops exponentially.
        \emph{Next steps:} compute the MPS ground state at $\chi \in \{2,4,8\}$
        via DMRG, load into a qiskit Statevector, rerun PITE and measure
        Trotter-step reduction; extend to 8- and 12-site TIM to see if the
        warm-start advantage grows with system size.

  \item \textbf{Chemistry benchmark vs.\ VQE and standard QITE.} How does PITE
        perform on H\textsubscript{2}, LiH, H\textsubscript{4} (JW-transformed
        minimal basis) versus VQE and measurement-based QITE, in terms of
        ground-state accuracy at fixed shot budget and two-qubit gate count?
        \emph{Basis:} the paper tests only TIM and simplified Hubbard; real
        molecular Hamiltonians have many more Pauli terms and larger
        coefficient dynamic range, so the practical claim is unquantified.
        \emph{Next steps:} build H\textsubscript{2}/STO-3G qubit Hamiltonian
        with OpenFermion; implement PITE, standard QITE (McArdle et al.\
        2019), and UCCSD-VQE in the same simulator; at matched shot budgets
        $\{10^5, 10^6, 10^7\}$ plot $|E - E_{\rm FCI}|$ vs.\ shot count for
        each method and report two-qubit gate counts.

  \item \textbf{Hardware-noise robustness.} How does the non-unitary Trotter
        step behave under realistic noise (mid-circuit measurement error on
        the ancilla, $T_1/T_2$ decoherence across the CNOT-ladder\,+\,$R_x$\,+\,
        measurement\,+\,reset sequence, residual system--ancilla entanglement
        after discard), and does noise translate into a benign $p_{\rm success}$
        reduction or into a systematic bias in the converged energy?
        \emph{Basis:} all results in the paper and this replication are ideal
        statevector; a 45-step run on 4-site TIM with $\sim 10$ Pauli terms
        involves $\sim 450$ mid-circuit measurements per shot, and at 99\%
        fidelity that is an $e^{-4.5}\!\sim\!1\%$ overhead on top of the
        $\sim 0.6\%$ algorithmic post-selection.
        \emph{Next steps:} run the ancilla-circuit form of PITE
        (\texttt{qiskit\_full\_ite.py}) under a calibrated depolarising
        noise model, sweep noise strength, separate coherent-gate from
        measurement error, and test whether readout-error mitigation or ZNE
        restore the ideal convergence.

  \item \textbf{Adaptive Trotter step size.} Does an adaptive schedule --- large
        $\Delta\tau$ early (when initial-state error dominates), small
        $\Delta\tau$ near convergence (when Trotter error dominates) --- reduce
        the total number of ancilla measurements to reach a target energy
        accuracy, and can it be learned online from observed per-step
        $p_{\rm success}$?
        \emph{Basis:} the paper uses a fixed $\Delta\tau = 0.1$; our data
        show the residual $|\Delta E|\sim 2\times 10^{-4}$ plateau is
        Trotter-limited, so shrinking $\Delta\tau$ late should improve it
        near-linearly, and $p_{\rm success}$ itself is a natural feedback
        signal.
        \emph{Next steps:} extend \texttt{ite\_tim.py} to accept a per-step
        $\Delta\tau$ schedule; compare Baseline (fixed 0.1) against Schedule~A
        (decreasing staircase) and Schedule~B (adaptive halving); test on the
        harder 4-site Hubbard regime the paper estimates at $\sim 10^9$
        shots and see whether the schedule fits inside $10^7$.
\end{enumerate}
